\section{Global stability from finite noisy observations}
\label{app:finite-noisy}

This section closes the observation-level gap between an exact-product
finite-probe reconstruction and a full-operator inverse bound.
It is an explicit measurement-level consequence of
Theorem~\ref{thm:response_stability} and
Corollary~\ref{cor:joint-stability}, not a new general inverse theorem.
The conclusion is global over all feasible factors in a specified
nondegenerate domain. No numerical recovery algorithm enters its proof.

\paragraph{Domain, observations, and quantifiers.}
Fix $L\ge K\ge2$, $D-K\ge L$, and known
$\eta_B,\eta_Z,\tau>0$, with $\lambda=\eta_Z/\tau^2$.
For positive $\alpha,s_A,s_B,M$, let
\[
 \mathcal D=\left\{(A,B):
 \begin{array}{l}
 A\in\mathbb R^{L\times K},\ A\mathbf1_K=\mathbf1_L,\
 A_{ij}\ge\alpha,\ \sigma_{\min}(A)\ge s_A,\\
 B\in\mathbb R^{K\times D},\
 \sigma_{\min}(B)\ge s_B,\ \|B\|_F\le M
 \end{array}\right\}.
\]
The margins describe an assumed parameter domain; they are not inferred
from one candidate's conditioning.
Given any observed product $\widehat\Theta\in\mathbb R^{L\times D}$,
choose an orthonormal-row matrix $U\in\mathbb R^{K\times D}$ from its top
$K$ right singular vectors, and set $\widehat P=U^\top U$.
Choose $W\in\mathbb R^{L\times D}$ with
$WW^\top=I_L$ and $WU^\top=0$; this is possible since $D-K\ge L$.
All choices depend only on $\widehat\Theta$ and dimensions, with a fixed
tie-breaking rule if needed. Writing $u_p^\top$ for row $p$ of $U$, define
\begin{equation}
 G_p=\frac{W+\mathbf1_Lu_p^\top}{\sqrt{2L}},
 \qquad p=1,\ldots,K.
 \label{eq:finite-noisy-probes}
\end{equation}
Orthogonality gives $\|G_p\|_F=1$. The complement code $W$ is deliberately
repeated in every query. Probes are selected before seeing their responses.

For any observed response matrices $\widehat R_p$, define
\begin{equation}
 \mathcal F=\left\{(A,B)\in\mathcal D:
 \begin{array}{l}
 \|AB-\widehat\Theta\|_F\le\varepsilon_\Theta,\\
 \bigl(\sum_{p=1}^K
 \|\mathcal R_{A,B}(G_p)-\widehat R_p\|_F^2\bigr)^{1/2}
 \le\varepsilon_R
 \end{array}\right\}.
 \label{eq:finite-noisy-feasible}
\end{equation}
This is the entire feasible set in the original $D$-dimensional basis
space, not a fit restricted to $\widehat P$.
Noise may be deterministic, adversarial, and dependent on the product
observation or probes, provided these norm bounds hold.

\paragraph{Explicit constants.}
Put
\begin{align}
 h&=\sqrt{2L}\bigl(K^{-1/2}+1+\sqrt L\bigr),&
 \beta&=\frac{8\lambda M}{s_A}(h\sqrt K+1),\\
 \kappa&=(s_B^{-2}+4s_A^{-2})^{1/2},&
 L_R&=\left[(2\eta_B\sqrt L+24\lambda M^2)^2
                         +(4\lambda M)^2\right]^{1/2}.
 \label{eq:finite-noisy-bridge-constants}
\end{align}
To make the inherited inverse constant fully explicit, define
\begin{align}
 L_M&=\max\{1,4\sqrt L/s_A\},&
 c_G&=\frac{4L}{\eta_Bs_A^2},&
 c_H&=\frac{8}{\lambda s_B^2},\\
 \mu_*&=\frac{\alpha}{2}
              \left(1+\frac{1}{4L_M^2}\right),&
 c_J&=\frac{L_M^2(c_G+c_H)}{\mu_*},&
 c_D&=\frac{2\sqrt{LK}}{s_A}c_J,\\
 c_{\rm row}
 &=\left[Kc_D^2+(c_G+Kc_D^2)^2\right]^{1/2},&
 \delta_0&=\min\left\{1,\frac1{2c_G},
                      \frac1{2\sqrt K\,c_{\rm row}}\right\}.
 \label{eq:finite-noisy-expanded}
\end{align}
Finally set
\begin{align}
 C_*&=\max\left\{
 2\sqrt{L+L_M^2M^2}\sqrt K\,c_{\rm row},
 \frac{4\sqrt{L+M^2}}{\delta_0}\right\},\\
 C_R&=2hC_*,\\
 C_\Theta&=\beta C_*+2\kappa(1+C_*L_R),\\
 \varepsilon_0&=\min\{\alpha s_B/4,s_As_B/8\}.
 \label{eq:finite-noisy-final-constants}
\end{align}
These constants can be very large. No sharpness or useful numerical
accuracy at a particular finite noise level is asserted.

\begin{theorem}[Whole-feasible-set diameter from $K$ noisy probes]
\label{thm:finite-noisy-diameter}
Under the preceding definitions, if
$0\le\varepsilon_\Theta\le\varepsilon_0$ and $\varepsilon_R\ge0$, then
\begin{equation}
 \operatorname{diam}_{d_{\rm perm}}(\mathcal F)
 \le C_\Theta\varepsilon_\Theta+C_R\varepsilon_R,
 \qquad
 d_{\rm perm}(x,x')=
 \min_{\Pi\in\mathfrak S_K}
 \left(\|A'-A\Pi\|_F^2+
       \|B'-\Pi^\top B\|_F^2\right)^{1/2}.
 \label{eq:finite-noisy-diameter}
\end{equation}
No smallness restriction on $\varepsilon_R$ is needed.
Empty feasible sets have diameter zero by convention; recovery
interpretations require a nonempty feasible set.
\end{theorem}

\begin{proof}
We give the observation-to-operator bridge first, then apply the explicit
global factor inverse to the original, unprojected factors.

\emph{Step 1: uniform control outside the observed subspace.}
For every $(A,B)\in\mathcal F$, $AB$ is a rank-$K$ competitor in the best
rank-$K$ approximation of $\widehat\Theta$. Hence
\[
 \|\widehat\Theta(I-\widehat P)\|_F\le\varepsilon_\Theta,
 \qquad
 \|AB(I-\widehat P)\|_F\le2\varepsilon_\Theta.
\]
Since $A$ has a left inverse with $\|A^\dagger\|_2\le s_A^{-1}$,
\begin{equation}
 \|B-B\widehat P\|_F\le2\varepsilon_\Theta/s_A.
 \label{eq:finite-noisy-tail}
\end{equation}
Set $\overline B=B\widehat P$. This is an analytical device, not an
assumption that the true basis lies in the estimated subspace.
For $S_i=\operatorname{diag}(a_i)-a_i a_i^\top$,
$\|S_i\|_2\le1$ and $\|\overline B\|_2,\|B\|_2\le M$, so
\[
 \|B^\top S_i^2B-\overline B^\top S_i^2\overline B\|_2
 \le2M\|B-\overline B\|_F.
\]
The shared Gram term is unchanged; the remaining operator difference
is block diagonal. Thus, uniformly over $\mathcal F$,
\begin{equation}
 \|\mathcal R_{A,B}-\mathcal R_{A,\overline B}\|_{F\to F}
 \le t,\qquad t=\frac{4\lambda M}{s_A}\varepsilon_\Theta.
 \label{eq:finite-noisy-compression}
\end{equation}
No rank property of $\overline B$ is required.

\emph{Step 2: stable decoding on a common support.}
Consider the larger linear class
$\mathcal S(G)=QG+\mathcal T(GU^\top)U$,
where $Q\in\mathbb R^{L\times L}$ and $\mathcal T$ right-multiplies row $i$
by an arbitrary $K\times K$ matrix $T_i$. Symmetry is unnecessary.
For the difference of two such operators, let
$E_p=\Delta\mathcal S(G_p)$ and
$E=(\sum_p\|E_p\|_F^2)^{1/2}$.
Set $c=(2L)^{-1/2}$. Complement projection gives
$E_pW^\top=c\,\Delta Q$, whence
\begin{equation}
 \|\Delta Q\|_F\le\frac{E}{c\sqrt K}.
 \label{eq:finite-noisy-Q}
\end{equation}
Projection onto $U$ gives, row by row,
\[
 E_pU^\top
 =c(\Delta Q\mathbf1_L)e_p^\top
   +c\,[e_p^\top\Delta T_i]_{i=1}^L.
\]
Stacking the $K$ equations enumerates every row of every $\Delta T_i$.
The triangle inequality and
$\|\Delta Q\mathbf1_L\|_2\le\sqrt L\,\|\Delta Q\|_F$ give
\begin{equation}
 \left(\sum_i\|\Delta T_i\|_F^2\right)^{1/2}
 \le E/c+\sqrt K\|\Delta Q\mathbf1_L\|_2
 \le (1+\sqrt L)E/c.
 \label{eq:finite-noisy-T}
\end{equation}
The shared term has induced norm $\|\Delta Q\|_2$ and the local
direct sum has norm at most $\max_i\|\Delta T_i\|_2$.
Equations~\ref{eq:finite-noisy-Q}--\ref{eq:finite-noisy-T} therefore imply
\begin{equation}
 \|\Delta\mathcal S\|_{F\to F}\le h E.
 \label{eq:finite-noisy-norming}
\end{equation}
For a projected factor response, take $Q=\eta_BAA^\top$ and
$T_i=\lambda U\overline B^\top S_i^2\overline B U^\top$.
The lemma applies even if those projected factors fail the domain's
rank margins.

\emph{Step 3: return to the original full operators.}
Choose any $x=(A,B),x'=(A',B')\in\mathcal F$.
Their observed response stacks differ by at most $2\varepsilon_R$.
Each unit-norm probe incurs at most $t$ error under
Equation~\ref{eq:finite-noisy-compression}; the stack incurs at most
$\sqrt K\,t$ for each pair. Applying
Equation~\ref{eq:finite-noisy-norming}, then adding the two full-operator
compression errors, yields
\begin{align}
 \|\mathcal R_x-\mathcal R_{x'}\|_{F\to F}
 &\le h(2\varepsilon_R+2\sqrt K\,t)+2t\\
 &=2h\varepsilon_R+\beta\varepsilon_\Theta.
 \label{eq:finite-noisy-full-response}
\end{align}
Every pair uses the same observed $U,W$ and probes.

\emph{Step 4: the global factor inverse.}
The two original factors obey the domain bounds with
$L_A=\sqrt L$ and $L_B=M$, and
$\|AB-A'B'\|_F\le2\varepsilon_\Theta$.
Our threshold ensures
$2\varepsilon_\Theta\le\min\{\alpha s_B/2,s_As_B/4\}$,
exactly the product-error condition of
Corollary~\ref{cor:joint-stability}.
That corollary's $C_*$ is the constant of
Theorem~\ref{thm:response_stability} evaluated at
$(\alpha/2,s_A/2,s_B/2,2\sqrt L,2M)$, which expands to
Equations~\ref{eq:finite-noisy-expanded}--\ref{eq:finite-noisy-final-constants}.
It includes the bounded-domain argument for large response differences,
not just a local inverse estimate.
Consequently,
\[
 d_{\rm perm}(x,x')
 \le C_*\|\mathcal R_x-\mathcal R_{x'}\|_{F\to F}
       +\kappa(1+C_*L_R)\|AB-A'B'\|_F.
\]
Substituting Equation~\ref{eq:finite-noisy-full-response} proves the claimed
bound for an arbitrary pair in $\mathcal F$, hence for its entire diameter.
No connected-component or candidate-neighborhood premise was used.
\end{proof}

\paragraph{Recovery interpretation and boundaries.}
If the true factors belong to $\mathcal D$ and satisfy the declared noise
bounds, they lie in $\mathcal F$, so every feasible factor estimate obeys
the displayed error bound. This does not prove that a particular solver
finds a feasible point. It does not verify domain membership from noisy
data, establish optimality of $K$ queries, or cover passive task-gradient
queries or unknown AdamW state. Each query returns a full $L\times D$
response, not one scalar.

No smooth choice of singular-vector basis is needed: the proof holds for
every chosen best rank-$K$ right subspace and orthonormal bases.
It does not assume a singular-value gap in the observation. For a nonempty
feasible set under the stated threshold, separation at the $K$th cut also
follows from the rank margins and perturbation inequalities, but that
fact is not needed in the proof.

For $K=1$, row stochasticity fixes $A=\mathbf1_L$ and the response's
softmax term vanishes. Directly,
$d_{\rm perm}(x,x')=\|B-B'\|_F\le2\varepsilon_\Theta/\sqrt L$.
Thus the boundary case needs no response information or small-noise
threshold. The nontrivial theorem concerns $K\ge2$.
